\documentclass[letterpaper,10pt]{article}

% ── Packages ──────────────────────────────────────────────────────────────────
\usepackage[left=0.55in, right=0.55in, top=0.5in, bottom=0.5in]{geometry}
\usepackage{lmodern}
\usepackage[T1]{fontenc}
\usepackage[utf8]{inputenc}
\usepackage{enumitem}
\usepackage[hidelinks]{hyperref}
\usepackage{titlesec}
\usepackage{parskip}

% ── Hyperlink colours (black = ATS safe) ──────────────────────────────────────
\hypersetup{
    colorlinks=true,
    linkcolor=black,
    urlcolor=black,
    citecolor=black
}

% ── Section formatting ────────────────────────────────────────────────────────
\titleformat{\section}
  {\large\bfseries\uppercase}
  {}{0em}{}
  [\titlerule]
\titlespacing{\section}{0pt}{6pt}{3pt}

% ── List spacing ──────────────────────────────────────────────────────────────
\setlist[itemize]{
  leftmargin=1.2em,
  itemsep=1.5pt,
  parsep=0pt,
  topsep=1pt
}

% ── No page numbers ───────────────────────────────────────────────────────────
\pagestyle{empty}

\begin{document}

% ── HEADER ────────────────────────────────────────────────────────────────────
\begin{center}
  {\LARGE \textbf{HERMAN PRAKOSO}} \\[3pt]
  \textit{Web Programmer} \textbf{|} \textit{PHP \& Basis Data} \textbf{|} \textit{System Analyst} \\[3pt]
  Temanggung, Jawa Tengah, Indonesia $\bullet$ +62 896-7364-7187 $\bullet$ \href{mailto:herman.prakoso33@gmail.com}{herman.prakoso33@gmail.com} \\
  Portofolio: \href{https://hermanprakoso.vercel.app}{hermanprakoso.vercel.app} $\bullet$ GitHub: \href{https://github.com/Tenku01}{github.com/Tenku01}
\end{center}

\vspace{1pt}

% ── RINGKASAN PROFESIONAL ─────────────────────────────────────────────────────
\section{Ringkasan Profesional}

Lulusan S1 Sistem Informasi Universitas Teknologi Yogyakarta (IPK 3.92/4.00, Cum Laude) dengan kompetensi kuat dalam dasar pemrograman web (PHP, JavaScript, HTML5/CSS3), perancangan basis data relasional MySQL, serta pengujian dan pelaporan sistem (SRS \& SDD). Berpengalaman membangun aplikasi web berbasis Laravel dan MySQL, merancang alur sistem informasi klinis/kesehatan (Juara 1 Lomba UI/UX), serta terbiasa menyelesaikan proyek perangkat lunak dari analisis kebutuhan hingga pengujian fungsional. Memiliki komitmen tinggi dan siap berkontribusi optimal 6 hari kerja per minggu sebagai peserta magang pemrograman klinik.

% ── PENDIDIKAN ────────────────────────────────────────────────────────────────
\section{Pendidikan}

\noindent
\textbf{Universitas Teknologi Yogyakarta} \hfill \textit{Lulusan 2026} \\
\textit{S1 Sistem Informasi} \textbf{|} \textbf{IPK: 3.92 / 4.00 (Cum Laude)} \hfill \textit{Yogyakarta, Indonesia}
\begin{itemize}
  \item \textbf{Fokus Bidang:} Rekayasa Perangkat Lunak, Pemrograman Web, Pemodelan Basis Data, Pengujian Sistem. Skripsi: Perancangan Sistem Informasi Web Multi-Role Menggunakan PHP Laravel, MySQL, dan Real-time WebSocket.
\end{itemize}

% ── KEAHLIAN TEKNIS ───────────────────────────────────────────────────────────
\section{Keahlian Teknis}

\begin{itemize}
  \item \textbf{Pemrograman \& Tampilan Web:} PHP, JavaScript, TypeScript, HTML5, CSS3, Tailwind CSS, Alpine.js, Laravel Framework, RESTful API.
  \item \textbf{Basis Data:} MySQL, Relational Database Modeling, Normalisasi Tabel, Query Optimization, Entity Relationship Diagram (ERD).
  \item \textbf{Pengujian \& Pelaporan Teknis:} Functional Testing, Black Box Testing, Bug Reporting, Software Requirements Specification (SRS), Software Design Description (SDD), User Manual.
  \item \textbf{Desain Antarmuka \& Alur Sistem:} Figma (UI/UX, Wireframe, Prototyping), System Workflow (Flowchart, DFD), UML/Use Case Modeling.
  \item \textbf{Tools \& Lingkungan Kerja:} Visual Studio Code, Git, GitHub, Postman, Draw.io.
\end{itemize}

% ── PENGALAMAN KERJA & MAGANG ─────────────────────────────────────────────────
\section{Pengalaman Kerja \& Magang}

\noindent
\textbf{Programmer \& System Analyst} \textbf{|} PT AKA Mendunia Sejahtera \hfill \textit{Feb 2026 -- Jun 2026} \\
\textit{Temanggung, Indonesia}
\begin{itemize}
  \item Mengembangkan aplikasi web sistem informasi menggunakan PHP (Laravel) dan MySQL dengan implementasi hak akses 5 role pengguna.
  \item Merancang struktur tabel basis data relasional, modul reservasi, serta integrasi payment gateway untuk otomasi verifikasi transaksi real-time.
  \item Mengimplementasikan fitur komunikasi real-time menggunakan WebSocket serta membangun antarmuka web interaktif berbasis JavaScript.
  \item Menyusun dokumen teknis (SRS \& SDD) dan menjalankan functional testing untuk memastikan keandalan alur transaksi sistem.
\end{itemize}

\noindent
\textbf{UI/UX Designer \& System Analyst Intern} \textbf{|} PT Swevel Universal Media \hfill \textit{Sep 2025 -- Jan 2026} \\
\textit{Yogyakarta, Indonesia}
\begin{itemize}
  \item Merancang arsitektur antarmuka web dan layout responsif di Figma dengan memprioritaskan kemudahan pengguna dan konsistensi visual.
  \item Menyusun dokumentasi kebutuhan perangkat lunak mencakup SRS, SDD, dan panduan penggunaan sistem (User Manual).
  \item Berkolaborasi dengan tim software engineer untuk memastikan akurasi implementasi tampilan web dan logika sistem.
\end{itemize}

\noindent
\textbf{UI/UX Designer \& System Analyst Intern} \textbf{|} PT Gama Integra Informatika \hfill \textit{Jan 2025 -- Jun 2025} \\
\textit{Yogyakarta, Indonesia}
\begin{itemize}
  \item Menganalisis kebutuhan fungsional sistem informasi dan alur bisnis operasional melalui wawancara stakeholder dan riset pengguna.
  \item Membangun wireframe, user flow, dan prototipe komponen web interaktif guna mendukung percepatan pengerjaan pengembang.
\end{itemize}

% ── PROYEK TERKAIT & PENGHARGAAN ──────────────────────────────────────────────
\section{Proyek Terkait \& Prestasi}

\noindent
\textbf{LIFE+ Health SuperApp} -- \textbf{Juara 1 Lomba UI/UX Design Tingkat Universitas} \hfill \textit{2025} \\
\textit{Teknologi: Figma, User Research, Information Architecture, Clinical Service Flow Design}
\begin{itemize}
  \item Merancang konsep antarmuka dan alur operasional sistem kesehatan yang mencakup modul reservasi layanan medis klinik, riwayat konsultasi pasien, dan pemantauan kesehatan berkala.
\end{itemize}

\noindent
\textbf{Portal Layanan \& Transparansi Keuangan (Website Dusun Demangan)} \hfill \textit{2026} \\
\textit{Teknologi: PHP, Laravel, Tailwind CSS, Alpine.js, MySQL, REST API}
\begin{itemize}
  \item Mengembangkan sistem web pelayanan publik terintegrasi mencakup pencatatan buku kas transaksi masuk/keluar, katalog \& approval peminjaman inventaris, serta laporan rekapitulasi iuran warga.
\end{itemize}

\end{document}

